\chapter{Interpretability and Model Governance}\label{ch:ml:interpretability-and-model-governance}

A risk committee asks why the model cut its position in a stock by half on Tuesday. The researcher's answer, that the
gradient-boosted model's output fell, is true and useless. The committee wants to know which inputs moved, whether the
model's response to them makes sense, whether it would have said the same last month, and who checked. This chapter is about
answering those questions honestly: explanations that are right about the model even when the features are correlated,
explanations that survive a retrain, and the documents (a \omterm{def:ml:interpretability-and-model-governance:card}{model card}, a validation report) through which a learned model
enters the firm's model risk management (Book~6, chapter~26). Two measurements anchor it. With two features correlated at
0.99, \omterm{def:ml:interpretability-and-model-governance:pd}{partial dependence} misstates a feature's effect by twice as much as \omterm{def:ml:interpretability-and-model-governance:pd}{accumulated local effects}. And on a problem with
a realistic signal-to-noise ratio, the top five features by \omterm{def:ml:feature-engineering-selection-and-importance:shapley}{SHAP value} are never the same five across ten retrains.

\section{What an explanation is for}

\begin{definition}[Interpretability]\label{def:ml:interpretability-and-model-governance:interp}
\emph{Interpretability}\index{interpretability} is the degree to which a person can understand why a model produces its
outputs: globally (how the prediction depends on each input across the data) or locally (why this prediction, for this
input). An explanation is itself a model of the model, and has an accuracy of its own.
\end{definition}

Explanations serve different readers. A researcher uses them to debug: a feature that matters and should not is a leak
(chapter~3). A risk manager uses them to check that the model's responses have the right sign and size, and that it will not
do something absurd in a market it has not seen. A committee or a regulator uses them to decide whether the firm understands
what it runs. Rudin (2019) argued that for high-stakes decisions an interpretable model should be preferred to an explained
black box; in trading, where the models are often small and the signal weak, that is a live option (chapter~4's baselines),
and a learned model should earn its complexity against it. The chapter's model is a gradient-boosted regression (chapter~5)
fitted to $y = x_1 + x_2^2 + 2x_3x_4 + \varepsilon$ on 6\,000 observations, where $x_2$ is, in two variants, a noisy copy of
$x_1$ (correlations 0.92 and 0.99).

\section{Global explanations: partial dependence and accumulated local effects}

\begin{definition}[Partial dependence, individual conditional expectation, accumulated local effects]\label{def:ml:interpretability-and-model-governance:pd}
The \emph{partial dependence}\index{partial dependence} of a model $f$ on feature $j$ is $\mathrm{PD}_j(v) = n^{-1}\sum_i
f(v, x_{i,-j})$: the average prediction with feature $j$ set to $v$ for every observation (Friedman, 2001). The
\emph{individual conditional expectation}\index{individual conditional expectation} (ICE) curves are the terms of that
average, one curve per observation (Goldstein and co-authors, 2015). \emph{Accumulated local effects}\index{accumulated local
effects} (ALE) average, within narrow intervals of feature $j$, the change in prediction when $j$ moves across the interval
for the observations that lie in it, and accumulate those changes (Apley and Zhu, 2020).
\end{definition}

\omterm{def:ml:interpretability-and-model-governance:pd}{Partial dependence} asks what the model predicts for combinations of features that may never occur: with $x_2$ nearly equal
to $x_1$, setting $x_1 = -0.9$ for an observation with $x_2 = 0.9$ asks the model about a point where it has no data, and a
tree ensemble answers with whatever its last split there says. ALE only moves each observation across a short interval around
its own value (\cref{lst:ml:gov:ale}). \Cref{tab:ml:gov:effects} measures both against the true effect of $x_1$, a straight
line of slope 1, and \cref{fig:ml:gov:pdale} shows them at correlation 0.99.

\begin{table}[htbp]
\centering\small
\begin{tabular}{@{}lccc@{}}
\toprule
& \multicolumn{3}{c}{correlation of $x_1$ and $x_2$}\\
\cmidrule(l){2-4}
root mean square error against the true effect of $x_1$ & 0 & 0.92 & 0.99\\
\midrule
\omterm{def:ml:interpretability-and-model-governance:pd}{partial dependence} & 0.027 & 0.134 & 0.229\\
\omterm{def:ml:interpretability-and-model-governance:pd}{accumulated local effects} & 0.028 & 0.044 & 0.110\\
\bottomrule
\end{tabular}
\caption{Error of two global explanations of the same fitted model, evaluated at the data's values of $x_1$ (the true effect
ranges over 2). Data: \texttt{ml\_explain.effects}.}\label{tab:ml:gov:effects}
\end{table}

\begin{omfigure}
\begin{tikzpicture}
\begin{axis}[width=10.5cm,height=6cm,xlabel={$x_1$},ylabel={centred effect on the prediction},tick label style={font=\small},
  label style={font=\small},xmin=-1,xmax=1,grid=major,grid style={gray!25},legend cell align=left,
  legend style={font=\footnotesize,at={(0.02,0.98)},anchor=north west,fill=white,draw=none},table/col sep=comma]
\addplot[black,thick,dashed] table[x=x,y=truth]{figdata/ml/21-interpretability-and-model-governance/pdale.csv};
\addplot[omThm,very thick,mark=*,mark size=1.3pt] table[x=x,y=pd]{figdata/ml/21-interpretability-and-model-governance/pdale.csv};
\addplot[omDef,very thick,mark=square*,mark size=1.3pt] table[x=x,y=ale]{figdata/ml/21-interpretability-and-model-governance/ale.csv};
\legend{true effect,partial dependence,accumulated local effects}
\end{axis}
\end{tikzpicture}

\omcaption{The effect of $x_1$ on a gradient-boosted model when $x_2$ is a copy of $x_1$ with correlation 0.99. Data:
\texttt{ml\_explain.curves}.}\label{fig:ml:gov:pdale}
\end{omfigure}

With independent features the two agree. As the correlation rises, \omterm{def:ml:interpretability-and-model-governance:pd}{partial dependence} absorbs part of $x_2$'s quadratic
effect into $x_1$ and its error grows five- and eightfold; ALE's grows too, because a model trained on nearly identical
features cannot itself say which one carries which effect, but it stays at half of \omterm{def:ml:interpretability-and-model-governance:pd}{partial dependence}'s. No explanation can
separate effects that the data do not separate; ALE at least does not add errors from places the data never visited.

ICE curves show what an average hides. The \omterm{def:ml:interpretability-and-model-governance:pd}{partial dependence} of $x_3$ is flat (slope 0.009): its effect is entirely its
interaction with $x_4$. The ICE curves of 30 observations (\cref{fig:ml:gov:ice}) fan out with slopes whose standard
deviation is 0.91 and whose correlation with the observation's $x_4$ is 0.98: the model has learned the interaction, and the
average of its responses has erased it.

\begin{omfigure}
\begin{tikzpicture}
\begin{axis}[width=10.5cm,height=6cm,xlabel={$x_3$},ylabel={prediction},tick label style={font=\small},
  label style={font=\small},xmin=-0.9,xmax=0.9,grid=major,grid style={gray!25},table/col sep=comma]
\pgfplotsinvokeforeach{0,1,2,3,4,5,6,7,8,9,10,11,12,13,14,15,16,17,18,19,20,21,22,23,24,25,26,27,28,29}{
  \addplot[omProp!45,thin] table[x=x,y=c#1]{figdata/ml/21-interpretability-and-model-governance/ice.csv};
}
\addplot[black,ultra thick] table[x=x,y=pd]{figdata/ml/21-interpretability-and-model-governance/ice.csv};
\end{axis}
\end{tikzpicture}

\omcaption{ICE curves of $x_3$ for 30 observations (thin) and their average (thick), for the model with independent features.
Data: \texttt{ml\_explain.ice\_x3}.}\label{fig:ml:gov:ice}
\end{omfigure}

\begin{definition}[Global surrogate]\label{def:ml:interpretability-and-model-governance:surrogate}
A \emph{global surrogate}\index{global surrogate} is an interpretable model (a shallow tree, a linear model) fitted to a
complex model's predictions; its \emph{fidelity}, the share of the predictions' variance it reproduces, says how much of the
model it explains.
\end{definition}

A depth-three tree fitted to the model's predictions splits only on $x_1$ and $x_2$ and reproduces 48\% of their variance: it
explains the additive half of the model and none of the interaction. A surrogate is an honest explanation only with its
fidelity printed next to it.

\section{Local explanations and their stability}

\begin{definition}[Counterfactual explanation]\label{def:ml:interpretability-and-model-governance:cf}
A \emph{counterfactual explanation}\index{counterfactual explanation} of a prediction is the smallest change of the inputs that
would have changed the decision, for example the value of one feature at which the forecast changes sign (Wachter,
Mittelstadt and Russell, 2017).
\end{definition}

Local explanations answer the committee's question about Tuesday. \omterm{def:ml:feature-engineering-selection-and-importance:shapley}{SHAP values} (chapter~6) split one prediction among the
features; a counterfactual says what would have had to be different. For one observation of the chapter's model, whose
prediction is 0.82 with $x_1 = 0.46$, the forecast turns negative only if $x_1$ falls to $-0.44$: most of this prediction
comes from the interaction, and $x_1$ alone would need a large move to reverse it.

An explanation that changes with every retrain explains the training run, not the market. \Cref{fig:ml:gov:ice} and the
table above come from one fit; the chapter also retrains a model ten times on bootstrap resamples of a problem with three
strong features (coefficients 1, 0.7, 0.5), five weak ones (0.2 each) and twelve with no effect, and ranks the features by
mean absolute \omterm{def:ml:feature-engineering-selection-and-importance:shapley}{SHAP value} each time (\cref{lst:ml:gov:stab}). When the features explain 66\% of the variance, the top five
features are the same set in 22\% of retrains, and the first retrain's top five keep their order with a mean Kendall
correlation of 0.89: the three strong features are always first, and which two weak features complete the list is a coin toss
among five equals. When they explain 5\%, closer to a return forecast, the top five are never the same set and the Kendall
correlation falls to 0.78; in some retrains a null feature enters the top five. A report that names ``the model's five most
important features'' from one fit is reporting noise in its last two places.

\section{Documentation: the model card}

\begin{definition}[Model card]\label{def:ml:interpretability-and-model-governance:card}
A \emph{model card}\index{model card} is a short document that travels with a model and states its purpose, training data,
evaluation, performance by relevant slices, behaviour, limitations and approvals (Mitchell and co-authors, 2019); in a firm
it is the learned model's entry in the model inventory (Book~6, chapter~26).
\end{definition}

The card is generated, not written: its fields are filled from the training run and the validation report, and the generator
refuses to call a card complete with fields missing. The chapter's card for its demonstration model lists fifteen fields and
reports two missing, monitoring and approvals, which is exactly the state of a model that has been built and validated but not
deployed. Its limitations field records the one fact a user must know: where $x_1$ and $x_2$ are correlated, read ALE, not
\omterm{def:ml:interpretability-and-model-governance:pd}{partial dependence}.

\begin{dated}{2026-09}{Model risk rules and learned models}\label{dat:ml:interpretability-and-model-governance:rules}
The United States' model risk guidance was revised on 17 April 2026 (SR 26-2, superseding SR 11-7); Book~6 (chapter~26)
describes both. The Prudential Regulation Authority's supervisory statement SS1/23, \emph{Model risk management principles for
banks}, published on 17 May 2023, sets five principles and applies them to all models used to inform business decisions,
``regardless of technology'', including ``the use of artificial intelligence in modelling techniques such as machine learning
to the extent that it applies to the use of models more generally''; its current version took effect on 23 April 2026,
after the PRA's April 2026 low-impact amendments.
\end{dated}

\section{Validating a learned model}

Validation of a learned model asks the questions of Book~6 (conceptual soundness, outcomes analysis, benchmarking,
effective challenge) plus a few that learned models make pressing: leakage, stability across retrains, and parity between the
model that was validated and the one that runs. \Cref{tab:ml:gov:checklist} is the chapter's checklist for its demonstration
model, filled where the evidence exists and marked ``not run'' where it does not; the numeric items run through Book~6's
\texttt{firm.modelval}.

\begin{table}[htbp]
\centering\small
\begin{tabular}{@{}p{4.3cm}lp{7.0cm}@{}}
\toprule
item & result & evidence\\
\midrule
out-of-sample performance & pass & $R^2$ 0.860 on 2\,000 new observations\\
challenger comparison & pass & ridge regression $R^2$ 0.349 on the same observations\\
implementation parity & pass & the exported trees (chapter~5) reproduce the model exactly on 200 rows\\
outcomes analysis & pass & 27 of 250 outcomes outside the 90\% interval (25 expected), binomial tail 0.37\\
conceptual soundness, leakage, stability, explanations, monitoring & not run & to be filled by the validator\\
\bottomrule
\end{tabular}
\caption{The validation checklist of the chapter's model. Data: \texttt{ml\_explain.validation}.}\label{tab:ml:gov:checklist}
\end{table}

\begin{method}[Explaining and governing a learned model]\label{met:ml:gov:method}
\begin{enumerate}
\item Explain effects with ALE where features are correlated, show ICE curves where interactions are suspected, and print a
surrogate's fidelity with the surrogate.
\item Report feature rankings from several retrains, with their stability; name only the features that are stable.
\item Answer questions about single decisions with \omterm{def:ml:feature-engineering-selection-and-importance:shapley}{SHAP values} and counterfactuals computed on the deployed model.
\item Generate the \omterm{def:ml:interpretability-and-model-governance:card}{model card} from the training and validation records; register it in the model inventory; let the validator
fill the checklist, and do not deploy with items not run.
\end{enumerate}
\end{method}

\section{Tutorial: explain it to the risk committee}\label{tut:ml:interpretability-and-model-governance:tut}

\begin{tutorial}
\textbf{Goal.} Explain a boosted model with PD, ICE, ALE and a surrogate; measure the stability of its SHAP rankings; generate
its card and checklist. \textbf{End state:} \cref{tab:ml:gov:effects,tab:ml:gov:checklist},
\cref{fig:ml:gov:pdale,fig:ml:gov:ice}.
\begin{tutsteps}
\item \textbf{\omterm{def:ml:interpretability-and-model-governance:pd}{Accumulated local effects}.}
\omcode{code/firm/modelcard/firm_modelcard.py}{50}{69}{First-order ALE with quantile bins.}\label{lst:ml:gov:ale}
\item \textbf{Stability of rankings across retrains.}
\omcode{code/firm/modelcard/firm_modelcard.py}{103}{112}{Top-set agreement and Kendall correlation across retrains.}\label{lst:ml:gov:stab}
\item \textbf{Run} \texttt{ml\_explain.effects()}, \texttt{ice\_x3()}, \texttt{surrogate()}, \texttt{stability(1.0)},
\texttt{stability(6.0)}, \texttt{validation()}, \texttt{card().render()} and \texttt{fig\_explain.py}.
\end{tutsteps}
\textbf{What to change next.} Add second-order ALE for the pair $(x_3, x_4)$; replace bootstrap retrains by retrains on
successive years; add a leakage canary (chapter~3) to the checklist.
\end{tutorial}

\section{Build: explanations and model cards}\label{bld:ml:interpretability-and-model-governance:modelcard}

\begin{build}
\bfield{Purpose} Explanations whose accuracy and stability are measured, and a learned model documented and validated like any
other model.
\bfield{Interface} \texttt{partial\_dependence}, \texttt{ice}, \texttt{ale}, \texttt{global\_surrogate},
\texttt{counterfactual}, \texttt{rank\_stability}; \texttt{ModelCard} with \texttt{render} and \texttt{missing};
\texttt{benchmark\_check}, \texttt{outcome\_check}, \texttt{validation\_checklist} (on \texttt{firm.modelval}).
\bfield{Rules} No explanation without its accuracy or fidelity; no ranking without its stability; no card with missing fields
marked complete; checklist items without evidence are ``not run''.
\bfield{Acceptance tests} \texttt{code/firm/modelcard/tests/}: PD and ALE recover a known additive function; ALE is right and
PD wrong on a correlated pair with an off-manifold trap; ICE slopes recover an interaction; a surrogate of a tree is perfect;
counterfactuals cross the threshold; rank stability is 1 for identical rankings; the card lists its missing fields.
\bfield{Stretch} Second-order ALE; SHAP interaction values; cards in the experiment tracker (chapter~25).
\end{build}

\begin{omsources}
\item J.~H.~Friedman, ``Greedy function approximation: a gradient boosting machine'', \emph{Annals of Statistics} 29(5), 2001.
\item A.~Goldstein, A.~Kapelner, J.~Bleich and E.~Pitkin, ``Peeking inside the black box: visualizing statistical learning with
plots of individual conditional expectation'', \emph{Journal of Computational and Graphical Statistics} 24(1), 2015.
\item D.~W.~Apley and J.~Zhu, ``Visualizing the effects of predictor variables in black box supervised learning models'',
\emph{Journal of the Royal Statistical Society B} 82(4), 2020.
\item M.~Mitchell and co-authors, ``Model cards for model reporting'', \emph{FAT*}, 2019.
\item S.~Wachter, B.~Mittelstadt and C.~Russell, ``Counterfactual explanations without opening the black box: automated
decisions and the GDPR'', SSRN 3063289, 2017.
\item C.~Rudin, ``Stop explaining black box machine learning models for high stakes decisions and use interpretable models
instead'', \emph{Nature Machine Intelligence} 1, 2019.
\item Prudential Regulation Authority, SS1/23, \emph{Model risk management principles for banks}, 2023 (current version
2026).
\end{omsources}

\section{Exercises}

\begin{exercise}[$\star$]\label{exo:ml:interpretability-and-model-governance:1}
For $f(x_1, x_2) = x_1 + x_2^2$ with $x_1$ and $x_2$ independent and uniform on $[-1, 1]$, write the \omterm{def:ml:interpretability-and-model-governance:pd}{partial dependence} on
$x_1$. What changes if $x_2 = x_1$ exactly and the model is $f$ itself?
\end{exercise}

\begin{exercise}[$\star$]\label{exo:ml:interpretability-and-model-governance:2}
With 250 outcomes and a 90\% interval, how many exceptions are expected? Is 27 alarming?
\end{exercise}

\begin{exercise}[$\star$]\label{exo:ml:interpretability-and-model-governance:3}
A \omterm{def:ml:interpretability-and-model-governance:surrogate}{global surrogate} has fidelity 0.48. What can and cannot be said from it?
\end{exercise}

\begin{exercise}[$\star\star$]\label{exo:ml:interpretability-and-model-governance:4}
Why is the \omterm{def:ml:interpretability-and-model-governance:pd}{partial dependence} of $x_3$ flat although $x_3$ matters?
\end{exercise}

\begin{exercise}[$\star\star$]\label{exo:ml:interpretability-and-model-governance:5}
Why do the fourth and fifth places of the SHAP ranking change across retrains even with a strong signal?
\end{exercise}

\begin{exercise}[$\star\star$]\label{exo:ml:interpretability-and-model-governance:6}
\emph{Find the flaw.} ``The model's \omterm{def:ml:feature-engineering-selection-and-importance:shapley}{SHAP values} show that the earnings-revision feature drives 40\% of the forecast, so the
model is a revisions model and the committee can think of it that way.''
\end{exercise}

\begin{exercise}[$\star\star\star$]\label{exo:ml:interpretability-and-model-governance:7}
\emph{Coding.} Rerun the stability experiment at the high signal level with ten retrains on disjoint tenths of the data
(400 observations each) instead of bootstrap resamples of all 4\,000. Does stability rise or fall, and why?
\end{exercise}

\begin{exercise}[$\star\star\star$]\label{exo:ml:interpretability-and-model-governance:8}
Show that for an additive model $f = \sum_jf_j(x_j)$ the ALE of feature $j$ equals $f_j$ up to a constant whatever the
dependence between the features, and that \omterm{def:ml:interpretability-and-model-governance:pd}{partial dependence} does too. Where, then, does \omterm{def:ml:interpretability-and-model-governance:pd}{partial dependence} go wrong?
\end{exercise}

\section{Problem: Explain It to the Risk Committee}

\begin{problem}[{Weekend problem --- an explanation has an error bar}]\label{pb:ml:interpretability-and-model-governance:1}
The chapter's model, retrains, card and checklist.

\medskip\noindent\textbf{Part I --- Global effects.}
\begin{enumerate}
\item What are the errors of PD and ALE at each correlation?
\item Why does ALE's error grow too?
\item What do the ICE curves of $x_3$ show that PD does not?
\item What does the surrogate explain, and what does it miss?
\end{enumerate}

\medskip\noindent\textbf{Part II --- Local and stable.}
\begin{enumerate}[resume]
\item What does the counterfactual say about the example's prediction?
\item How stable is the top five at each signal level?
\item What should a report on feature importance contain?
\item How would you answer the committee's question about Tuesday?
\end{enumerate}

\medskip\noindent\textbf{Part III --- Governance.}
\begin{enumerate}[resume]
\item What does the \omterm{def:ml:interpretability-and-model-governance:card}{model card} contain, and what is missing?
\item Which checklist items pass, on what evidence?
\item Why is implementation parity on the list?
\item What does the dated box say about how supervisors treat learned models?
\end{enumerate}

\medskip\noindent\textbf{Part IV --- The verdict.}
\begin{enumerate}[resume]
\item State the \emph{named result}: \omterm{def:ml:interpretability-and-model-governance:pd}{partial dependence}'s error against the true effect under correlated features compared with
ALE's, and the rank stability of the top five SHAP features across ten retrains.
\item Would you replace the boosted model by an interpretable one here?
\item Which tier (Book~6) would you give a learned alpha model, and why?
\item What does effective challenge look like for a learned model?
\item What must be re-validated after a retrain?
\item How do you explain a model whose inputs are embeddings (chapter~10)?
\item What should never be in a \omterm{def:ml:interpretability-and-model-governance:card}{model card}?
\item In one sentence: what is an explanation?
\end{enumerate}
\end{problem}

\section{Interview questions}

\begin{interviewq}[$\star$ \iqroles{researcher, mle}]\label{iq:ml:interpretability-and-model-governance:1}
What is the difference between \omterm{def:ml:interpretability-and-model-governance:pd}{partial dependence} and ALE, and when does it matter?
\end{interviewq}

\begin{interviewq}[$\star\star$ \iqroles{researcher}]\label{iq:ml:interpretability-and-model-governance:2}
How stable are SHAP-based feature rankings, and how would you measure it?
\end{interviewq}

\begin{interviewq}[$\star\star$ \iqroles{mle}]\label{iq:ml:interpretability-and-model-governance:3}
What goes into a \omterm{def:ml:interpretability-and-model-governance:card}{model card} for a trading model?
\end{interviewq}

\begin{interviewq}[$\star\star$ \iqroles{researcher, trader}]\label{iq:ml:interpretability-and-model-governance:4}
A portfolio manager asks why the model sold a stock. What do you show?
\end{interviewq}

\begin{interviewq}[$\star\star$ \iqroles{mle, researcher}]\label{iq:ml:interpretability-and-model-governance:5}
How would you validate a gradient-boosted model before production?
\end{interviewq}

\begin{interviewq}[$\star\star\star$ \iqroles{researcher}]\label{iq:ml:interpretability-and-model-governance:6}
When would you prefer an interpretable model to an explained black box in trading?
\end{interviewq}
